% =====================================================================
\chapter{Cuestiones}
\label{cap:test}
% =====================================================================

Veinte preguntas de opción múltiple, de 1,5 puntos cada una (30 puntos en total). Cada pregunta tiene una sola respuesta correcta; las respuestas erróneas no restan. Tiempo orientativo: 45 minutos.

\subsection*{Módulo 20. Ajuste fino}

\pregunta{En el ajuste supervisado (SFT) de un modelo de chat, la entropía cruzada se calcula normalmente:}
{sobre todos los tokens de la secuencia, incluidos los de la instrucción.}
{solo sobre los tokens de la respuesta del asistente, incluido el token de fin de turno que la cierra.}
{solo sobre el último token de la respuesta.}
{solo sobre los tokens de la instrucción, que son los que escribe el usuario.}

\pregunta{LoRA de rango $r=8$ aplicado a una matriz $W\in\R^{4096\times 1024}$ entrena:}
{$20\,480$ parámetros.}
{$32\,768$ parámetros.}
{$40\,960$ parámetros.}
{$4\,194\,304$ parámetros.}

\pregunta{En LoRA se inicializa $B=0$ y $A$ con entradas aleatorias porque:}
{así la matriz $BA$ es invertible desde el primer paso.}
{se reduce a la mitad la memoria del estado del optimizador.}
{así $A$ no recibe gradiente durante todo el entrenamiento.}
{el modelo ajustado empieza siendo exactamente el preentrenado y, aun así, $B$ recibe un gradiente no nulo en el primer paso.}

\pregunta{En el modelo cuadrático del olvido, $\mathcal{L}(\theta)=(\theta-a)^{2}+\lambda(\theta-\theta_0)^{2}$, con $\theta_0=0$, $a=4$ y $\lambda=3$, el minimizador $\theta^{\ast}$ vale:}
{$1$}
{$4/3\approx 1{,}33$}
{$3$}
{$4$}

\pregunta{Una cuantización por bloques usa índices de $b=4$ bits y una constante de escala en coma flotante de 32 bits por cada bloque de $32$ pesos. ¿Cuántos bits ocupa cada parámetro?}
{$4{,}5$}
{$5$}
{$8$}
{$36$}

\subsection*{Módulo 21. Aprendizaje por refuerzo y preferencias}

\pregunta{Un episodio produce las recompensas $R_1=1$, $R_2=0$ y $R_3=2$ y termina. Con factor de descuento $\gamma=0{,}5$, el retorno $G_0$ vale:}
{$0{,}5$}
{$1$}
{$1{,}25$}
{$1{,}5$}

\pregunta{Restar una línea base $b(s)$, que no depende de la acción, al retorno en el estimador del gradiente de política:}
{introduce un sesgo que hay que corregir con muestreo de importancia.}
{no cambia la esperanza del estimador y puede reducir su varianza.}
{anula el gradiente cuando $b(s)$ es la recompensa media.}
{solo es válido si $b(s)=0$ para todos los estados.}

\pregunta{En el objetivo recortado de PPO con $\varepsilon=0{,}2$, un término con cociente $\rho=0{,}7$ y ventaja $A=-1$ vale:}
{$0{,}7$}
{$-0{,}7$}
{$-0{,}8$}
{$-1{,}2$}

\pregunta{En 10 comparaciones entre dos respuestas, la respuesta~$a$ se prefiere 8 veces. Bajo el modelo de Bradley--Terry, la estimación de máxima verosimilitud de $r_a-r_b$ es:}
{$0{,}8$}
{$\ln 1{,}25\approx 0{,}223$}
{$4$}
{$\ln 4\approx 1{,}386$}

\pregunta{Al comenzar el entrenamiento con DPO, con $\pi_\theta=\piref$, la pérdida de cada par de preferencias vale:}
{$0$}
{$1$}
{$\beta$}
{$\ln 2\approx 0{,}693$}

\pregunta{En GRPO, si las $G$ respuestas a una misma instrucción obtienen exactamente la misma recompensa:}
{todas sus ventajas son nulas y el grupo no aporta gradiente al término recortado.}
{todas sus ventajas valen~1 y el grupo se refuerza por igual.}
{la división por la desviación típica produce ventajas infinitas, sea cual sea $\varepsilon_0$.}
{desaparece la penalización KL del objetivo.}

\subsection*{Módulo 22. Agentes, recuperación y RAG}

\pregunta{Una colección tiene $N=1000$ documentos y el término $t$ aparece en $10$ de ellos. Con el logaritmo natural, como en los apuntes, $\idf_t$ vale:}
{$2$}
{$4{,}61$}
{$6{,}91$}
{$100$}

\pregunta{En BM25 con $k_1=1{,}5$, cuando la frecuencia $f$ del término en el documento tiende a infinito, el factor de saturación $g(f,\ell)$ tiende a:}
{$1$}
{$1{,}5$}
{$2{,}5$}
{$+\infty$}

\pregunta{Un ranking tiene las relevancias (N, R, N, R, R), donde R indica relevante, y hay $\abs{\Rel}=4$ documentos relevantes en la colección. Su precisión media (AP) es:}
{$0{,}4$}
{$0{,}5$}
{$0{,}533$}
{$0{,}6$}

\pregunta{Se dice que un \emph{benchmark} está contaminado cuando:}
{sus preguntas o sus respuestas aparecen en los datos de entrenamiento del modelo evaluado.}
{contiene preguntas ambiguas o con varias respuestas válidas.}
{se evalúa con un sistema RAG en lugar de con el modelo solo.}
{mide la latencia en lugar de la exactitud.}

\pregunta{Un agente encadena 20 pasos y cada uno sale bien, de forma independiente, con probabilidad $0{,}98$. La probabilidad de completar la tarea es, aproximadamente:}
{$0{,}400$}
{$0{,}600$}
{$0{,}668$}
{$0{,}980$}

\subsection*{Módulo 23. Interpretabilidad y seguridad}

\pregunta{La hipótesis de superposición afirma que una red neuronal:}
{dedica exactamente una neurona a cada concepto.}
{no tiene flujo residual en las capas intermedias.}
{solo puede aprender funciones lineales de su entrada.}
{representa más características que dimensiones, como direcciones casi ortogonales, lo que es rentable cuando las características son dispersas.}

\pregunta{En un autoencoder disperso con penalización $L^{1}$, el código óptimo de un latente aislado es $f^{\ast}=\max(a-\lambda/2,\,0)$. Con preactivación $a=0{,}3$ y $\lambda=0{,}8$, $f^{\ast}$ vale:}
{$0$}
{$-0{,}1$}
{$0{,}3$}
{$0{,}7$}

\pregunta{El escalado de temperatura de un clasificador, $\softmax(z/T)$ con $T>0$ ajustado en validación:}
{cambia su exactitud, porque puede alterar la clase de mayor probabilidad.}
{cambia solo las confianzas y deja intacta la clase predicha.}
{cambia el número de clases efectivas del modelo.}
{invierte el orden de los logits cuando $T>1$.}

\pregunta{En 600 intentos de \emph{red teaming} independientes no se observa ningún éxito del ataque. Una cota superior aproximada, con confianza del $95\,\%$, para su probabilidad de éxito por intento es:}
{$0\,\%$}
{$0{,}05\,\%$}
{$0{,}5\,\%$}
{$3\,\%$}
